\documentclass[11pt]{article}

\usepackage[margin=0.86in]{geometry}
\usepackage{amsmath,amssymb,amsthm}
\usepackage{graphicx}
\usepackage{microtype}
\usepackage[hidelinks]{hyperref}

\newcommand{\R}{\mathbb R}
\newcommand{\1}{\mathbf 1}
\newcommand{\Sp}{S_p}
\newtheorem{theorem}{Theorem}

\title{An analytic counterexample to removing transversality\\
from the local geometry of Schur idempotents}
\author{FA/Banach run \texttt{fa\_banach\_001}, lane 3}
\date{August 13, 2026}

\begin{document}
\maketitle

\section*{Source metadata and status}

The source is J.~Parcet, M.~de la Salle, and E.~Tablate,
\emph{The local geometry of idempotent Schur multipliers},
arXiv:2312.02895v2, published in \emph{Forum of Mathematics, Pi} 13 (2025),
e14 \cite{PST}.  Section 1.7 of the arXiv PDF (Section 2.7 in the published
version) asks whether the equivalence in Theorem A survives without
transversality when the second component of the boundary normal is nonzero.
The authors expect counterexamples, but suggest that analytic boundaries
probably do not admit them.

\medskip
\noindent
\textbf{Claimed status: full counterexample, likely valid.}
Transversality is necessary even for a polynomial hypersurface in
$\R^2\times\R^2$.

\begin{center}
  \includegraphics[width=0.98\textwidth]{figures/open_problem_crop.png}
\end{center}

\section*{Proof intuition}

Make many boundary sections pass through one degenerate fiber while their
tangent directions vary.  Ordinarily such a continuum of directions would
trigger the paper's Meyer--Fefferman obstruction.  Here every meeting point
has vanishing first normal component, so that blow-up cannot see it.

Boundedness comes from a different feature.  The sign of one $y$-coordinate
splits the symbol into just two pieces.  On one piece it is the triangular
projection, and on the other it is the same projection with the order
reversed.  Thus a changing analytic family of boundary lines is hidden inside
a finite measurable pasting of familiar bounded multipliers.

\section*{Counterexample theorem}

\begin{theorem}\label{thm:main}
Let $U=V=(-1,1)^2$, write $x=(r,q)\in U$ and $y=(s,t)\in V$, and put
\[
  \Sigma=\{(x,y)\in U\times V:s>r^2t\}.
\]
For every $1<p<\infty$, the indicator $\1_\Sigma$ is a bounded Schur
multiplier on
\[
  \Sp\bigl(L_2(U),L_2(V)\bigr).
\]
At $z_0=((0,0),(0,0))\in\partial\Sigma$, the $V$-component of the
boundary normal is nonzero, but the zero-curvature condition (2) and the
triangular representation (3) of Theorem A both fail on every neighborhood
of $z_0$.  The boundary is real analytic.
\end{theorem}

\section*{Proof}

Define
\[
  F(r,q,s,t)=s-r^2t.
\]
Then $\Sigma=\{F>0\}$ and
\[
  \nabla F=(-2rt,0,1,-r^2).
\]
In particular, $\nabla F$ never vanishes, so $\partial\Sigma$ is a
real-analytic hypersurface.  At $z_0$ its two normal components are
\[
  \mathbf n_1(z_0)=(0,0),\qquad
  \mathbf n_2(z_0)=(1,0).
\]
Thus $z_0$ is nontransverse but satisfies precisely the weaker normal
hypothesis proposed in the source question.

We next prove boundedness.  Let
\[
  V_+=\{(s,t)\in V:t>0\},\qquad
  V_-=\{(s,t)\in V:t<0\},
\]
and define measurable scalar maps
\[
  a(r,q)=r^2,\qquad b(s,t)=\frac{s}{t}\quad(t\ne0).
\]
If $\tau_<(u,v)=\1_{\{u<v\}}$ and
$\tau_>(u,v)=\1_{\{u>v\}}$, then, outside the null set $t=0$,
\begin{equation}\label{eq:split}
  \1_\Sigma(x,y)
  =\1_{V_+}(y)\,\tau_<\bigl(a(x),b(y)\bigr)
   +\1_{V_-}(y)\,\tau_>\bigl(a(x),b(y)\bigr).
\end{equation}
The symbols $\tau_<$ and $\tau_>$ are the continuous triangular and
reverse-triangular Schur projections.  They are bounded on $\Sp$ for every
$1<p<\infty$; denote a common bound by $C_p$.

The pushforward of Lebesgue measure on $U$ under $a$ is absolutely continuous
with respect to Lebesgue measure on $\R$ (on $(0,1)$ its density is a constant
multiple of $u^{-1/2}$).  The pushforwards of Lebesgue measure on $V_+$ and
$V_-$ under $b$ are also absolutely continuous: for a null set $A\subset\R$,
Fubini's theorem gives
\[
  \bigl|\{(s,t)\in V_\pm:s/t\in A\}\bigr|=0.
\]
The source's change-of-variables lemma \cite[Lemma 1.1]{PST} therefore gives
\[
 \bigl\|\tau_<\circ(a\times b)\bigr\|_{MS_p(U,V_+)}\le C_p,
 \qquad
 \bigl\|\tau_>\circ(a\times b)\bigr\|_{MS_p(U,V_-)}\le C_p.
\]
Restriction to $V_+$ or $V_-$ is left multiplication by an orthogonal
projection on $L_2(V)$ and hence contractive on every Schatten class.
Equation \eqref{eq:split} now yields
\[
  \|\1_\Sigma\|_{MS_p(U,V)}\le 2C_p<\infty.
\]

It remains to check that the other two conditions fail.  For fixed
$x=(r,q)$,
\[
  \partial\Sigma_x=\{(s,t):s=r^2t\},
  \qquad
  T_{(0,0)}\partial\Sigma_x
    =\operatorname{span}\{(r^2,1)\}.
\]
Every section passes through $y_0=(0,0)$, but these tangent lines vary with
$r$.  Choosing $r=0$ and any sufficiently small $r\ne0$ disproves the
zero-curvature condition on every product neighborhood of $z_0$.

For completeness, suppose there were neighborhoods $U_0,V_0$ and
$C^1$ functions $f_1:U_0\to\R$, $f_2:V_0\to\R$ such that
\[
  \Sigma\cap(U_0\times V_0)
  =\{(x,y):f_1(x)>f_2(y)\}.
\]
Since $(x,y_0)$ is a boundary point for every $x\in U_0$, continuity forces
$f_1(x)=f_2(y_0)$ for every such $x$.  Hence the set on the right is
independent of $x$.  This is impossible: choose small $\rho,t>0$ with
$x^0=(0,0),x^\rho=(\rho,0)\in U_0$ and
$y=(\rho^2t/2,t)\in V_0$.  Then $(x^0,y)\in\Sigma$ while
$(x^\rho,y)\notin\Sigma$.  Thus condition (3) also fails. \qed

\section*{Scope, novelty check, and review recommendation}

The example gives a negative answer to both parts of the source question:
the equivalence fails without transversality, and it can fail for an analytic
boundary.  Both factor manifolds have dimension two.  The result makes no
claim about the paper's separate global questions for Fourier multipliers on
$\mathrm{SL}_2(\R)$.

Run-index and bounded web searches through August 13, 2026 used the arXiv id,
the exact transversality paragraph, analytic-boundary variants, and the
formula $s>x^2t$.  They found the 2025 publication and a 2025 thesis still
presenting the transverse theorem, but no later solution or matching
counterexample.  This is a bounded novelty audit, not a priority claim.

\medskip
\noindent
\textbf{Human review recommendation.} Review as a likely valid full
counterexample.  The key audit points are the two-piece identity
\eqref{eq:split}, applicability of the source's pullback lemma to $a$ and
$b$, and the boundary-point argument excluding a single triangular model.

\begin{thebibliography}{9}
\bibitem{PST}
J.~Parcet, M.~de la Salle, and E.~Tablate,
\emph{The local geometry of idempotent Schur multipliers},
arXiv:2312.02895v2; Forum Math. Pi \textbf{13} (2025), e14,
\href{https://doi.org/10.1017/fmp.2025.6}{doi:10.1017/fmp.2025.6}.
\end{thebibliography}

\end{document}
